\documentclass[a4paper]{article}
\usepackage{geometry}
\usepackage{amssymb}
\usepackage{pdfpages}
\usepackage{graphbox}
\usepackage{tikz}
\begin{document}

\begin{center}
  {\Large\bf Percentiles of the Student $\mathbf{t}$ distribution}
\end{center}


\pagestyle{empty}
\thispagestyle{empty}
\centering
\begin{tabular}{r|ccccc}
  \begin{minipage}{20mm}degrees\\of freedom\end{minipage}
  & \begin{minipage}{1.3cm}\hfill{\Large $t_{0.9\hphantom{9}}$}\hfill\end{minipage}
  & \begin{minipage}{1.3cm}\hfill{\Large $t_{0.95}             $}\hfill\end{minipage}
  & \begin{minipage}{1.3cm}\hfill{\Large $t_{0.975}            $}\hfill\end{minipage}
  & \begin{minipage}{1.3cm}\hfill{\Large $t_{0.99\hphantom{9}} $}\hfill\end{minipage}
  & \begin{minipage}{1.3cm}\hfill{\Large $t_{0.995}            $}\hfill\end{minipage}
  \rule[-4mm]{0mm}{10mm}
  \\  \hline
\input student_table_values.txt  % created by student_number_maker.R
  \hline
\end{tabular}

\begin{centering}
\includegraphics[align=c,width=2.2in]{student_diagram.pdf}
\end{centering}

\noindent For example, to calculate the~$95^\mathrm{th}$ percentile on
one degree of freedom, the first line gives the different quantiles.
Choose the entry with column heading~$t_{0.95}$ to obtain~6.31.

\begin{tikzpicture}[remember picture,overlay]
  \node[anchor=south east,
    xshift=-1cm,
    yshift=1.1cm]
  at (current page.south east){\includegraphics[width=2in]{mono-logo.pdf}};
  \node[anchor=south west,
    xshift=1cm,
    yshift=1.1cm]
  at (current page.south west){{\tt https://github.com/RobinHankin/tables.git}};
\end{tikzpicture}


\end{document}
